\question{10}{Considere un grafo representado por las siguientes listas de adyacencia:

\begin{center}
\begin{tikzpicture}[baseline=(current bounding box.center),
  idxcell/.style={draw, minimum size=0.42cm, inner sep=0pt, outer sep=0pt},
  slot/.style={draw, minimum size=0.42cm, inner sep=1pt, outer sep=0pt, font=\scriptsize},
  -latex, thick,
]
  \node[idxcell]                    (idx0) at (0, 0) {};
  \node[idxcell, below=0pt of idx0] (idx1) {};
  \node[idxcell, below=0pt of idx1] (idx2) {};
  \node[idxcell, below=0pt of idx2] (idx3) {};
  \node[idxcell, below=0pt of idx3] (idx4) {};
  \node[idxcell, below=0pt of idx4] (idx5) {};

  \node[left=6pt of idx0] {$1$};
  \node[left=6pt of idx1] {$2$};
  \node[left=6pt of idx2] {$3$};
  \node[left=6pt of idx3] {$4$};
  \node[left=6pt of idx4] {$5$};
  \node[left=6pt of idx5] {$6$};

  % idx0 (v1): 3, 5
  \node[slot, right=0.35cm of idx0] (r0a1) {3};
  \node[slot, right=0pt of r0a1]    (r0a2) {};
  \node[slot, right=0.35cm of r0a2] (r0b1) {5};
  \node[slot, right=0pt of r0b1]    (r0b2) {};
  \draw (idx0.center) -- (r0a1);
  \draw (r0a2.center) -- (r0b1);

  % idx1 (v2): 3, 6
  \node[slot, right=0.35cm of idx1] (r1a1) {3};
  \node[slot, right=0pt of r1a1]    (r1a2) {};
  \node[slot, right=0.35cm of r1a2] (r1b1) {6};
  \node[slot, right=0pt of r1b1]    (r1b2) {};
  \draw (idx1.center) -- (r1a1);
  \draw (r1a2.center) -- (r1b1);

  % idx2 (v3): 1, 2, 4
  \node[slot, right=0.35cm of idx2] (r2a1) {1};
  \node[slot, right=0pt of r2a1]    (r2a2) {};
  \node[slot, right=0.35cm of r2a2] (r2b1) {2};
  \node[slot, right=0pt of r2b1]    (r2b2) {};
  \node[slot, right=0.35cm of r2b2] (r2c1) {4};
  \node[slot, right=0pt of r2c1]    (r2c2) {};
  \draw (idx2.center) -- (r2a1);
  \draw (r2a2.center) -- (r2b1);
  \draw (r2b2.center) -- (r2c1);

  % idx3 (v4): 3, 5
  \node[slot, right=0.35cm of idx3] (r3a1) {3};
  \node[slot, right=0pt of r3a1]    (r3a2) {};
  \node[slot, right=0.35cm of r3a2] (r3b1) {5};
  \node[slot, right=0pt of r3b1]    (r3b2) {};
  \draw (idx3.center) -- (r3a1);
  \draw (r3a2.center) -- (r3b1);

  % idx4 (v5): 1, 4, 6
  \node[slot, right=0.35cm of idx4] (r4a1) {1};
  \node[slot, right=0pt of r4a1]    (r4a2) {};
  \node[slot, right=0.35cm of r4a2] (r4b1) {4};
  \node[slot, right=0pt of r4b1]    (r4b2) {};
  \node[slot, right=0.35cm of r4b2] (r4c1) {6};
  \node[slot, right=0pt of r4c1]    (r4c2) {};
  \draw (idx4.center) -- (r4a1);
  \draw (r4a2.center) -- (r4b1);
  \draw (r4b2.center) -- (r4c1);

  % idx5 (v6): 2, 5
  \node[slot, right=0.35cm of idx5] (r5a1) {2};
  \node[slot, right=0pt of r5a1]    (r5a2) {};
  \node[slot, right=0.35cm of r5a2] (r5b1) {5};
  \node[slot, right=0pt of r5b1]    (r5b2) {};
  \draw (idx5.center) -- (r5a1);
  \draw (r5a2.center) -- (r5b1);
\end{tikzpicture}
\end{center}

El primer elemento del arreglo de apuntadores representa el vértice \(v_1\). Cada vez que visite los vecinos de un vértice, hágalo en el orden en que aparecen en su lista.
}

\begin{enumerate}
  \item (2 décimas) Dibuje el grafo que representan estas listas.
  \item (1 décimas) ¿El grafo es dirigido?
  \item (1 décimas) ¿El grafo es ponderado?
  \item (6 décimas) Ejecute DFS recursivo desde \(v_1\). Para cada paso (cada vez que visite un vértice por primera vez), muestre la pila de llamadas recursivas en ese momento y el arreglo de predecesores actualizado.
\end{enumerate}

\bigskip
